\documentclass[../NDCNNAI_main.tex]{subfiles}
\graphicspath{{\subfix{../fig/}}}

\begin{document}
\subsection{Теорема о представлении или почему архитектура с УАС почти всегда сеть прямого распространения}
В прикладных архитектурах часто присутствует <<преобразование представления>>: мы переводим объект в удобное пространство,
где стандартные сети эффективны, аппроксимируем там, и затем возвращаемся обратно (encoder/feature map $\to$ MLP $\to$ decoder).

Теорема о представлении Крациоса формализует, что подобная логика не случайна: любая универсальная архитектура
на широком классе пространств может рассматриваться как <<MLP-универсальность, перенесённая через подходящее преобразование>>.
\begin{definition}[Вложение (embedding)]\label{def:embedding_repr_style}
Отображение $\Psi:Y\to Z$ называется вложением (embedding), если оно непрерывно, инъективно и является
гомеоморфизмом на свой образ $\Psi(Y)$.
\end{definition}

\begin{theorem}[О представлении, Крациос \cite{Kratsios2021}]
Пусть $\sigma\colon\mathbb{R}\to\mathbb{R}$ непрерывна.
Обозначим через $(\mathcal{F}_0,\circ_0)$ стандартную архитектуру полносвязных нейронных сетей (MLP)
на $C(\mathbb{R})$ с активацией $\sigma$.
Пусть $X$ --- функциональное пространство, гомеоморфное бесконечномерному Фреше-пространству.
Если архитектура $(\mathcal{F},\circ)$ имеет УАС в $X$, то существует семейство вложений
\[
\{\Phi_i\colon C(\mathbb{R})\hookrightarrow X\}_{i\in I},
\]
такое что для любых $f\in X$ и $\varepsilon>0$ найдутся индекс $i\in I$, функция
$g_\varepsilon\in \mathcal{NN}(\mathcal{F}_0,\circ_0)$ и реализация
$f_\varepsilon\in \mathcal{NN}(\mathcal{F},\circ)$, удовлетворяющие
\begin{align}
d_X\bigl(f,\Phi_i(g_\varepsilon)\bigr) &< \varepsilon, \label{eq:repr_push_style}\\
d_{\mathrm{ucc}}\bigl(g_\varepsilon,\Phi_i^{-1}(f_\varepsilon)\bigr) &< \varepsilon. \label{eq:repr_pull_style}
\end{align}
\end{theorem}

\begin{proof}
\noindent\textbf{Шаг 1. Универсальность MLP на $C(\mathbb{R})$.}
Для непрерывной неполиномиальной активации $\sigma$ класс MLP плотен в $C(\mathbb{R})$ в топологии ucc.
Этот факт используется как стандартная теорема универсальной аппроксимации.

\noindent\textbf{Шаг 2. Характеризация.}
Из Леммы существует семейство $X_i\subset X$ и гомеоморфизмы
$\Phi_i:C(\mathbb{R})\to X_i$ такие, что $\bigcup_i X_i$ плотно в $X$ и
$\mathcal{NN}(\mathcal{F},\circ)\cap X_i$ плотно в $X_i$ для каждого $i$.

\noindent\textbf{Шаг 3. Доказательство \eqref{eq:repr_push_style} (push-forward).}
Зафиксируем $f\in X$ и $\varepsilon>0$.
По плотности $\bigcup_i X_i$ найдутся $i\in I$ и $f_i\in X_i$ такие, что
\begin{equation}\label{eq:repr_step1}
d_X(f,f_i)<\varepsilon/2.
\end{equation}
Так как $\Phi_i$ -- гомеоморфизм, он переносит плотность:
из плотности $\mathcal{NN}(\mathcal{F}_0,\circ_0)$ в $C(\mathbb{R})$ следует, что
$\Phi_i(\mathcal{NN}(\mathcal{F}_0,\circ_0))$ плотно в $X_i$.
Значит, существует $g_\varepsilon\in \mathcal{NN}(\mathcal{F}_0,\circ_0)$ такое, что
\begin{equation}\label{eq:repr_step2}
d_X\bigl(f_i,\Phi_i(g_\varepsilon)\bigr)<\varepsilon/2.
\end{equation}
Из \eqref{eq:repr_step1}--\eqref{eq:repr_step2} и неравенства треугольника получаем
\[
d_X\bigl(f,\Phi_i(g_\varepsilon)\bigr)\le d_X(f,f_i)+d_X\bigl(f_i,\Phi_i(g_\varepsilon)\bigr)<\varepsilon,
\]
то есть \eqref{eq:repr_push_style}.

\noindent\textbf{Шаг 4. Доказательство \eqref{eq:repr_pull_style} (pull-back).}
По Лемме множество $\mathcal{NN}(\mathcal{F},\circ)\cap X_i$ плотно в $X_i$.
Применяя гомеоморфизм $\Phi_i^{-1}\colon X_i\to C(\mathbb{R})$, получаем, что
\[
\Phi_i^{-1}\bigl(\mathcal{NN}(\mathcal{F},\circ)\cap X_i\bigr)
\quad\text{плотно в } C(\mathbb{R}) \text{ (в ucc)}.
\]
Следовательно, существует $\tilde g\in \Phi_i^{-1}\bigl(\mathcal{NN}(\mathcal{F},\circ)\cap X_i\bigr)$ такое, что
\[
d_{\mathrm{ucc}}(g_\varepsilon,\tilde g)<\varepsilon.
\]
По определению множества существует $f_\varepsilon\in \mathcal{NN}(\mathcal{F},\circ)\cap X_i$ с
$\tilde g=\Phi_i^{-1}(f_\varepsilon)$, откуда получаем \eqref{eq:repr_pull_style}.
\end{proof}

\begin{corollary}
Формулы \eqref{eq:repr_push_style}--\eqref{eq:repr_pull_style} читаются как <<двусторонний перенос>>:
\begin{itemize}
\item push-forward: чтобы приблизить $f\in X$, мы выбираем подходящее <<окно>> $X_i$ и приближаем $f$
элементами вида $\Phi_i(g)$, где $g$ -- MLP-функция на $C(\mathbb{R})$;
\item pull-back: любая универсальная архитектура на $X$ в каждом $X_i$ индуцирует плотный класс
в $C(\mathbb{R})$ после обратного перехода через $\Phi_i^{-1}$.
\end{itemize}
Тем самым результат теоермы даёт структурную классификацию:
все УАС-архитектуры на <<хороших>> пространствах эквивалентны (с точностью до вложений)
универсальности стандартных нейросетей на $C(\mathbb{R})$.    
\end{corollary}

\subsection{Теорема существования: <<малые>> универсальные аппроксиматоры}
Напомним ряд определений и фактов, которые были введены или использованы ранее.
\begin{definition}[Топологическая транзитивность (для отображения)]
Пусть $(X,\tau)$ --- топологическое пространство, $T\colon X\to X$ -- непрерывное отображение.
Говорят, что $T$ топологически транзитивно, если для любых непустых открытых множеств
$U,V\subset X$ найдётся $n\in\mathbb{N}$ такое, что
\[
T^{n}(U)\cap V\neq\varnothing.
\]
\end{definition}

\begin{definition}[Плотная орбита / гиперциклический элемент]
Для $x\in X$ орбита равна $\mathrm{Orb}_T(x)\coloneq \{T^n(x):n\in\mathbb{N}\}$.
Элемент $x$ называется гиперциклическим (или универсальным для динамики $T$),
если $\overline{\mathrm{Orb}_T(x)}=X$.
\end{definition}

\begin{theorem}[Биркгофа о транзитивности]
На <<хороших>> бэровских пространствах (в частности, сепарабельных Фреше-пространствах)
из топологическая транзитивность динамики следует существование множества (которое даже образует плотное $G_\delta$)
точек с плотной орбитой. В линейной динамике это формулируется как транзитивность непрерывного линейного оператора эквивалентна гиперцикличности.
\end{theorem}

\begin{remark}
Транзитивность означает <<динамика умеет перепрыгивать между любыми двумя открытыми областями>>.
В бэровских пространствах это почти автоматически <<склеивается>> в существование одной точки,
орбита которой заходит во все элементы счётной базы топологии -- отсюда плотность орбиты.
\end{remark}

\begin{definition}[Полугруппа операторов и орбита]
Пусть $(X,\tau)$ --- топологическое пространство, $\{T_\theta:X\to X\}_{\theta\in\Theta}$ --
семейство непрерывных отображений. Обозначим через $\mathsf{S}$ полугруппу,
порождённую $\{T_\theta\}$ относительно композиции.
Для $f\in X$ орбита равна
\[
\mathrm{Orb}_{\mathsf{S}}(f)\coloneq \{S(f):S\in\mathsf{S}\}\subset X.
\]
\end{definition}

\begin{definition}[Топологическая транзитивность действия полугруппы]
Действие $\mathsf{S}\curvearrowright X$ называется топологически транзитивным, если
для любых непустых открытых $U,V\subset X$ найдётся $S\in\mathsf{S}$ такое, что
$S(U)\cap V\neq\varnothing$.
\end{definition}

\begin{theorem}[О существовании, Крациос \cite{Kratsios2021}]
Пусть $(X,\tau)$ --- сепарабельное бэровское пространство.
Если действие полугруппы $\mathsf{S}\curvearrowright X$ топологически транзитивно, то существует
$f_\star\in X$ такое, что
\[
\overline{\mathrm{Orb}_{\mathsf{S}}(f_\star)}=X.
\]
В терминах архитектур: класс может иметь УАС даже при <<крошечном>> наборе примитивов
(в пределе только одного), если правило комбинирования способно породить транзитивную динамику.
\end{theorem}

\begin{proof}
Пусть $\{V_m\}_{m\in\mathbb{N}}$ --- счётная база топологии $\tau$.
Для каждого $m$ положим
\[
    U_m\coloneq \bigcup_{S\in\mathsf{S}} S^{-1}(V_m)=\{f\in X:\exists S\in\mathsf{S}\ \text{такое, что}\ S(f)\in V_m\}.
\]
Тогда $U_m$ открыто (объединение прообразов открытого множества под непрерывными $S$).

Покажем, что $U_m$ плотно: пусть $U\subset X$ непусто и открыто. По транзитивности
найдётся $S\in\mathsf{S}$ с $S(U)\cap V_m\neq\varnothing$, то есть
$U\cap S^{-1}(V_m)\neq\varnothing$, значит $U\cap U_m\neq\varnothing$.

Теперь применяем теорему Бэра о категории: пересечение $\bigcap_{m\in\mathbb{N}}U_m$ плотно (и, в частности, непусто).
Берём $f_\star\in\bigcap_m U_m$. Тогда для каждого $m$ существует $S_m\in\mathsf{S}$ с $S_m(f_\star)\in V_m$,
то есть орбита $\mathrm{Orb}_{\mathsf{S}}(f_\star)$ пересекает каждый элемент базы, а значит плотна в $X$.
\end{proof}

\begin{definition}[Отмеченное метрическое пространство и липшицевы отображения]
Отмеченное метрическое пространство --- тройка $(X,d_X,0_X)$, где $0_X\in X$ фиксирована.
Обозначим через $\mathrm{Lip}_0(X,\mathbb{R})$ множество липшицевых функций $f\colon X\to\mathbb{R}$,
удовлетворяющих $f(0_X)=0$.
\end{definition}

\begin{definition}[Свободное пространство (Lipschitz-free space) $\mathcal{B}(X)$ и отображение $\delta_X$]
Рассмотрим формальные конечные комбинации $\sum_{k=1}^N \alpha_k \delta_{x_k}$, где $x_k\in X$, $\alpha_k\in\mathbb{R}$.
Норму задаём по формуле
\[
\Bigl\|\sum_{k=1}^N \alpha_k \delta_{x_k}\Bigr\|_{\mathcal{B}(X)}
\coloneq \sup\Bigl\{\sum_{k=1}^N \alpha_k f(x_k)\ :\ f\in \mathrm{Lip}_0(X,\mathbb{R}),\ \mathrm{Lip}(f)\le 1\Bigr\}.
\]
Пополнение даёт банахово пространство $\mathcal{B}(X)$, а отображение $\delta_X\colon x\mapsto \delta_x$
является (нелинейной) изометрией $X\hookrightarrow \mathcal{B}(X)$.
\end{definition}

\begin{definition}[Барицентр и барицентричность]
Непрерывное отображение $\rho\colon \mathcal{B}\to X$ называется барицентром, если
$\rho\circ \delta_X = \mathrm{id}_X$. Если такой $\rho$ существует, то $X$ называется барицентрическим.
\end{definition}

\begin{theorem}[О существовании малых универсальных аппроксиматоров]
Пусть тройка $(X,d_X,0_X)$ --- сепарабельное отмеченное метрическое пространство, содержащее хотя бы две точки. $(\mathcal{X},d_{\mathcal{X}},0_{\mathcal{X}})$ -- функциональное пространство, рассматриваемое как отмеченное метрическое,
и $\mathcal{X}_0\subset \mathcal{X}$ --- плотное барицентрическое подпространство.
Тогда существует непустое множество индексов $I$ с предпорядком $\le$,
подмножество $\{x_i\}_{i\in I}\subset X\setminus\{0_X\}$ и тройки
$\{(B_i,\iota_i,\phi_i)\}_{i\in I}$, где $B_i$ --- линейные подпространства $\mathcal{B}(\mathcal{X}_0)$,
$\iota_i\colon \mathcal{B}(X)\to B_i$ -- ограниченные линейные изоморфизмы,
$\phi_i\colon \mathcal{B}(X)\to \mathcal{B}(X)$ -- ограниченные линейные отображения, такие что:
\begin{enumerate}[label=(\roman*)]
\item $\mathcal{B}(\mathcal{X}_0)=\bigcup_{i\in I} B_i$;
\item если $i\le j$, то $B_i\subseteq B_j$;
\item для каждого $i\in I$ множество
\[
\Bigl\{\ \iota_i\bigl(\phi_i^{\,n}(\delta_X(x_i))\bigr)\ :\ n\in\mathbb{N}_+\ \Bigr\}
\quad\text{плотно в }B_i\text{ (в относительной топологии);} 
\]
\item[(iv)] архитектура $(F,\circ)$ с примитивами $F=\{x_i\}_{i\in I}$ может быть определена так,
чтобы её множество реализаций содержало (по крайней мере) семейство
\[
\Bigl\{\ \rho\bigl(\iota_i(\phi_i^{\,n}(\delta_X(x_i)))\bigr)\ :\ i\in I,\ n\in\mathbb{N}_+\ \Bigr\}\subset \mathcal{X}_0,
\]
и тогда $(F,\circ)$ обладает УАС на $\mathcal{X}$, т.\,е.
$\overline{\mathcal{NN}(F,\circ)}^{\,\mathcal{X}}=\mathcal{X}$.
\end{enumerate}

Кроме того, если $\mathcal{X}=\mathcal{X}_0$, то можно выбрать $I$ одноточечным и $\iota$ тождественным
на $\mathcal{B}(\mathcal{X})$.
\end{theorem}

\begin{proof}
Зафиксируем плотное барицентрическое $\mathcal{X}_0\subset\mathcal{X}$ и рассмотрим $B(\mathcal{X}_0)$.

\noindent\textbf{Шаг 1: линейная динамика в $\mathcal{B}(X)$.}
Поскольку $X$ сепарабельно, то $\mathcal{B}(X)$ сепарабельно.
Для каждого $i$ выбирают $x_i\in X\setminus\{0_X\}$ и строят ограниченный линейный оператор
$\phi_i\colon \mathcal{B}(X)\to B(X)$ такой, что орбита $\{\phi_i^{\,n}(\delta_X(x_i))\}_{n\in\mathbb{N}_+}$ плотна в $B(X)$.

\noindent\textbf{Шаг 2: перенос плотности в подпространства $B_i\subset B(\mathcal{X}_0)$.}
Далее, пользуясь структурой $\mathcal{B}(\mathcal{X}_0)$ как (индуктивного) объединения банаховых подпространств,
строят направленную систему $\{B_i\}_{i\in I}$ и изоморфизмы $\iota_i\colon \mathcal{B}(X)\to B_i$;
тогда образы плотных орбит остаются плотными:
$\{\iota_i(\phi_i^{\,n}(\delta_X(x_i)))\}_n$ плотно в $B_i$.

\noindent\textbf{Шаг 3: склейка по $i\in I$ и возвращение в $\mathcal{X}_0$.}
Из свойств топологии на объединении $\bigcup_i B_i$ получают плотность множества
$\{\iota_i(\phi_i^{\,n}(\delta_X(x_i))) : i\in I,\ n\in\mathbb{N}_+\}$ в $\mathcal{B}(\mathcal{X}_0)$
(в тексте это оформлено через промежуточное пространство $\widetilde{B}$ и проверку того,
что любое непустое открытое множество пересекается с указанным семейством).

Поскольку $\mathcal{X}_0$ барицентрично, существует непрерывный барицентр
$\rho\colon\mathcal{B}(\mathcal{X}_0)\to\mathcal{X}_0$, причём $\rho\circ \delta_{\mathcal{X}_0}=\mathrm{id}_{\mathcal{X}_0}$.
Образ плотного множества под непрерывной сюръекцией плотен, поэтому
\[
\Bigl\{\ \rho(\iota_i(\phi_i^{\,n}(\delta_X(x_i))))\ :\ i\in I,\ n\in\mathbb{N}_+\ \Bigr\}
\]
плотно в $\mathcal{X}_0$.
Наконец, так как $\mathcal{X}_0$ плотно в $\mathcal{X}$, получаем плотность и в $\mathcal{X}$.
\end{proof}

\begin{remark}
Теорема утверждает, что <<универсальность>> можно получить,
если у нас есть (i) небольшой набор стартовых примитивов $F$ (вплоть до одного при $\mathcal{X}=\mathcal{X}_0$)
и (ii) правило $\circ$, способное реализовать достаточно богатую динамику (итерации операторов $\phi_i$)
после линеаризации через $\delta_X$ и обратного проектирования через барицентр $\rho$.
Это согласуется с лекционной идеей <<УАС как плотная достижимость>>.
\end{remark}

\subsection*{Глубокие сети прямого распространения через транзитивные композиционные операторы}
Рассмотрим стандартную feed-forward архитектуру: чередование аффинных преобразований и
покомпонентной нелинейности. Пусть $\sigma\colon \mathbb{R}\to\mathbb{R}$ --- непрерывная активация,
а $\sigma^\bullet\colon \mathbb{R}^m\to\mathbb{R}^m$ --- покомпонентное продолжение,
\[
\sigma^\bullet(x_1,\dots,x_m) \coloneq  (\sigma(x_1),\dots,\sigma(x_m)).
\]
Тогда типичный слой имеет вид $x\mapsto \sigma^\bullet(Ax+b)$.

\begin{definition}[Композиционный оператор]
Зафиксируем отображение входного пространства
\[
S_{A,b}(x)\coloneq \sigma^\bullet(Ax+b),\qquad A\in\mathbb{R}^{m\times m},\; b\in\mathbb{R}^m,
\]
и определим оператор на пространстве функций
\begin{equation}
\Phi_{A,b}\colon X_0\to X_0,\qquad (\Phi_{A,b}f)(x)\coloneq f(S_{A,b}(x))=f\circ S_{A,b}(x).
\label{eq:phi_def}
\end{equation}
Оператор $\Phi_{A,b}$ является линейным. При естественных условиях на $S_{A,b}$ он непрерывен в ucc-топологии;
в частности, нужная непрерывность обеспечивается в ключевом для приложений случае $A=I_m$.
\end{definition}

\begin{remark}
Если базовая сеть реализует некоторый аппроксиматор $f$, то добавление глубины за счёт повторения одного и того же
блока $S_{A,b}$ приводит к семейству
\[
f,\; f\circ S_{A,b},\; f\circ S_{A,b}^{2},\;\dots,\; f\circ S_{A,b}^{N},
\]
то есть к орбите $\{ \Phi_{A,b}^N(f)\}_{N\ge 0}$ в пространстве функций. Именно плотность таких орбит и связывается с УАС.
\end{remark}

По поводу топологической динамики на пространствах функций можно сделать следующее замечание.
\begin{remark}
Если оператор на пространстве функций обладает транзитивностью,
то итерации этого оператора способны <<разнести>> аппроксиматоры по всему пространству.
В терминах Крациоса: глубина (итерация) + транзитивность $\Rightarrow$ универсальность.
\end{remark}

Отметим, что на практике удобен следующее \textbf{достаточное условие топологической транзитивности оператора композиции} (в ucc-топологии): достаточно, чтобы функция активации не имела неподвижных точек и была инъективной.
В статье показано, что на самом деле это \textbf{критерий}.
Фиксируем основной для приложения случай
\[
A=I_m,\qquad b\in\mathbb{R}^m,\quad b_i>0,
\qquad
S(x)=\sigma^\bullet(x+b),
\qquad
\Phi f=f\circ S.
\]

\begin{theorem}[Глубина как транзитивная динамика]
Пусть $\sigma\colon \mathbb{R}\to\mathbb{R}$ непрерывна, а $b\in\mathbb{R}^m$ имеет строго положительные компоненты.
Тогда для оператора $\Phi\colon X_0\to X_0$, $\Phi(f)=f\circ\sigma^\bullet(\cdot+b)$, эквивалентны:
\begin{enumerate}
\item $\Phi$ топологически транзитивен на $X_0$ (в ucc-топологии);
\item $\Phi$ гиперцикличен на $X_0$;
\item $\sigma$ инъективна и не имеет неподвижных точек: $\forall x\in\mathbb{R}\;\sigma(x)\neq x$;
\item выполняется строгая альтернатива
\[
\forall x\in\mathbb{R}\;\sigma(x)>x
\qquad \text{или} \qquad
\forall x\in\mathbb{R}\;\sigma(x)<x.
\]
\end{enumerate}
\end{theorem}

\begin{remark}
Для непрерывной функции $h(x)\coloneq \sigma(x)-x$ отсутствие нулей означает, что $h$ не меняет знак
(иначе по теореме о промежуточных значениях существовал бы $x_\ast$ с $h(x_\ast)=0$).
Отсюда следует либо $\sigma(x)>x$ для всех $x$, либо $\sigma(x)<x$ для всех $x$.
\end{remark}

Ограничимся лишь доказательством леммы.
\begin{lemma}[Критерий транзитивности композиционного оператора]
Пусть $S\colon \mathbb{R}^m\to\mathbb{R}^m$ --- непрерывное инъективное отображение и
$\Phi\colon C(\mathbb{R}^m,\mathbb{R}^n)\to C(\mathbb{R}^m,\mathbb{R}^n)$ задано как $\Phi(f)=f\circ S$.
Тогда $\Phi$ топологически транзитивен тогда и только тогда, когда $S$ обладает свойством убегания компактов:
для любого компакта $K\subset\mathbb{R}^m$ существует $N\in\mathbb{N}$ такое, что
\begin{equation}
S^{N}(K)\cap K=\varnothing.
\label{eq:runaway}
\end{equation}
\end{lemma}

\begin{proof}[Схема доказательства]
\leavevmode
\begin{enumerate}
\item \underline{Необходимость:} если $\Phi$ транзитивен, то символ $S$ обязан быть инъективен и удовлетворять \eqref{eq:runaway},
иначе можно построить два непустых открытых множества в $C(\mathbb{R}^m,\mathbb{R}^n)$, которые не пересекаются
под действием итераций $\Phi^N$.
\item \underline{Достаточность:} если $S$ инъективен и выполняется \eqref{eq:runaway}, то для любых непустых открытых $U,V$
можно подобрать функцию $f\in U$ так, чтобы $f\circ S^N$ попала в $V$ для некоторого $N$.
Грубо говоря, условие \eqref{eq:runaway} позволяет независимо <<настроить>> значения функции на $K$ и на $S^N(K)$,
не создавая конфликтов на пересечении (его нет), после чего используется плотность элементарных модификаций
в ucc-топологии.
\end{enumerate}
Тем самым транзитивность $\Phi$ сводится к геометрическому свойству итераций $S$ на $\mathbb{R}^m$.
\end{proof}

\begin{remark}
Для $S(x)=\sigma^\bullet(x+b)$ с $b_i>0$ условие \eqref{eq:runaway} эквивалентно отсутствию неподвижных точек у $\sigma$:
если $\sigma(x)>x$ для всех $x$, то итерации $S^N$ монотонно уводят любой компакт <<в сторону>> (вдоль каждой координаты),
и пересечение с исходным компактом исчезает при большом $N$; аналогично для $\sigma(x)<x$.  
\end{remark}

Как добавление транзитивных динамических блоков может обеспечить УАС проясняет следующая пара утверждений.
\begin{proposition}[Механизм local-to-global]
Пусть $X_0=C(\mathbb{R}^m,\mathbb{R}^n)$ с ucc-топологией и пусть $\Phi\colon X_0\to X_0$ --- непрерывное отображение,
которое топологически транзитивно. Тогда для любого непустого открытого множества $U\subset X_0$ объединение
итерационных образов
\[
\bigcup_{N\ge 0}\Phi^N(U)
\]
плотно в $X_0$.
\end{proposition}

\begin{corollary}[Добавление глубины $\Rightarrow$ УАС на всём $X_0$]
Пусть архитектура $(F,\circ)$ реализует множество $\mathcal{N}\coloneq \mathcal{NN}(F,\circ)\subset X_0$,
и существует непустое открытое множество $U\subset X_0$ такое, что $\mathcal{N}$ плотно в $U$.
Если $\sigma$ транзитивна (в смысле Theorem 5), так что композиционный оператор $\Phi(f)=f\circ\sigma^\bullet(\cdot+b)$ транзитивен,
то углублённая архитектура, допускающая композиции вида
\[
f\circ \underbrace{\sigma^\bullet(\cdot+b)\circ \cdots \circ \sigma^\bullet(\cdot+b)}_{N\ \text{раз}},
\qquad f\in\mathcal{N},\ N\in\mathbb{N},
\]
обладает УАС на всём $X_0$.    
\end{corollary}
\begin{proof}
Пусть $V\subset X_0$ --- непустое открытое множество. Так как $\Phi$ транзитивен, по Proposition 1 найдутся
$N$ и $g\in U$ такие, что $\Phi^N(g)\in V$.
Так как $\mathcal{N}$ плотно в $U$, можно выбрать $f\in \mathcal{N}$ достаточно близко к $g$ (в ucc),
и по непрерывности $\Phi^N$ получим $\Phi^N(f)\in V$.
Но $\Phi^N(f)=f\circ(\sigma^\bullet(\cdot+b))^N$ реализуется углублённой архитектурой, значит $V$ пересекается
с множеством реализаций. Поскольку $V$ произвольно, это и есть плотность, то есть УАС.
\end{proof}
\end{document}